\section{Conclusion}
\label{sec:06_con}

This paper presented BiDecT, an encoder-free video captioning framework that integrates pre-trained multimodal representations directly into a bidirectional decoding architecture. The consistent state-of-the-art results across three benchmarks of varying scale and annotation density, combined with substantial computational savings over encoder-based counterparts, support the premise that intermediate feature encoding is unnecessary when sufficiently expressive pre-trained representations are available. These findings suggest that encoder-free bidirectional decoding is a robust and generalizable paradigm for video captioning.

A remaining limitation, identified through qualitative analysis, is the model’s difficulty in discriminating between visually similar fine-grained actions when the distinguishing cues lie beyond the scope of the current visual and semantic features. Integrating complementary modalities such as audio, or incorporating explicit action recognition signals, represents a promising direction for future work. More broadly, the framework may extend naturally to related vision-language tasks such as dense video captioning, video question answering, and multimodal summarization, where computational efficiency and linguistic coherence are jointly important.

